% SPDX-FileCopyrightText: 2024-2026 Fabian Januszewski
% SPDX-License-Identifier: CC-BY-4.0

%======================================================================
% gl2_preamble.tex --- shared preamble and macro file
%
%   The SINGLE source of macros for the unified contract (gl2_contract.tex)
%   and for every companion document that cites it.  Documents \input this
%   file; they do NOT copy it.  (The pre-2026-09 corpus copied macros between
%   hecke_action.tex, hecke_extensions.tex and monodromy_v5.tex; that is what
%   let three incompatible meanings of \De, \tau and \iota coexist.)
%
%   Usage, after \documentclass and before \begin{document}:
%       % SPDX-FileCopyrightText: 2024-2026 Fabian Januszewski
% SPDX-License-Identifier: CC-BY-4.0

%======================================================================
% gl2_preamble.tex --- shared preamble and macro file
%
%   The SINGLE source of macros for the unified contract (gl2_contract.tex)
%   and for every companion document that cites it.  Documents \input this
%   file; they do NOT copy it.  (The pre-2026-09 corpus copied macros between
%   hecke_action.tex, hecke_extensions.tex and monodromy_v5.tex; that is what
%   let three incompatible meanings of \De, \tau and \iota coexist.)
%
%   Usage, after \documentclass and before \begin{document}:
%       % SPDX-FileCopyrightText: 2024-2026 Fabian Januszewski
% SPDX-License-Identifier: CC-BY-4.0

%======================================================================
% gl2_preamble.tex --- shared preamble and macro file
%
%   The SINGLE source of macros for the unified contract (gl2_contract.tex)
%   and for every companion document that cites it.  Documents \input this
%   file; they do NOT copy it.  (The pre-2026-09 corpus copied macros between
%   hecke_action.tex, hecke_extensions.tex and monodromy_v5.tex; that is what
%   let three incompatible meanings of \De, \tau and \iota coexist.)
%
%   Usage, after \documentclass and before \begin{document}:
%       % SPDX-FileCopyrightText: 2024-2026 Fabian Januszewski
% SPDX-License-Identifier: CC-BY-4.0

%======================================================================
% gl2_preamble.tex --- shared preamble and macro file
%
%   The SINGLE source of macros for the unified contract (gl2_contract.tex)
%   and for every companion document that cites it.  Documents \input this
%   file; they do NOT copy it.  (The pre-2026-09 corpus copied macros between
%   hecke_action.tex, hecke_extensions.tex and monodromy_v5.tex; that is what
%   let three incompatible meanings of \De, \tau and \iota coexist.)
%
%   Usage, after \documentclass and before \begin{document}:
%       \input{gl2_preamble}
%
%   Every macro below is normative: the notation chapter of the contract
%   (gl2_notation.tex) is the prose form of this file, and the two must agree.
%   A symbol that is not declared here is not part of the fixed notation.
%
%   RETIRED MACROS.  Symbols retired by the notation chapter are redefined
%   here to raise a LaTeX error rather than silently typeset, so that text
%   re-homed from the old documents cannot carry an old meaning in unnoticed.
%   See the block "Retired symbols" at the end.
%======================================================================

\usepackage{amsmath, amssymb, amsthm}
\usepackage{mathtools}
\usepackage{geometry}
\usepackage{booktabs}
\usepackage{array}
\usepackage{longtable}
\usepackage{enumitem}
\usepackage{tikz}
\usepackage{tikz-cd}
\usepackage[colorlinks=true,linkcolor=blue,citecolor=blue,urlcolor=blue]{hyperref}
\setlist{nosep}
\emergencystretch=3em

%----------------------------------------------------------------------
% Theorem environments.  One counter, sectioned, so every statement has a
% unique number of the form <section>.<n> that can be cited from code.
%----------------------------------------------------------------------
\newtheorem{theorem}{Theorem}[section]
\newtheorem{lemma}[theorem]{Lemma}
\newtheorem{proposition}[theorem]{Proposition}
\newtheorem{corollary}[theorem]{Corollary}
\newtheorem{claim}[theorem]{Claim}

\theoremstyle{definition}
\newtheorem{definition}[theorem]{Definition}
\newtheorem{algorithm}[theorem]{Algorithm}
\newtheorem{obligation}[theorem]{Obligation}
\newtheorem{certificate}[theorem]{Certificate}

\theoremstyle{remark}
\newtheorem{remark}[theorem]{Remark}
\newtheorem{assumption}[theorem]{Assumption}
\newtheorem{example}[theorem]{Example}
\newtheorem{prediction}[theorem]{Prediction}
\newtheorem{openitem}[theorem]{Open construction item}

%----------------------------------------------------------------------
% Normative boxes.
%----------------------------------------------------------------------
\newenvironment{contractbox}
  {\par\medskip\noindent\begin{quote}\textbf{CONTRACT.}\ }
  {\end{quote}\medskip}
\newenvironment{decisionbox}
  {\par\medskip\noindent\begin{quote}\textbf{OWNER DECISION.}\ }
  {\end{quote}\medskip}
\newenvironment{erratumbox}
  {\par\medskip\noindent\begin{quote}\textbf{ERRATUM.}\ }
  {\end{quote}\medskip}

%----------------------------------------------------------------------
% STATUS tags.  Every non-definitional assertion carries exactly one.
%----------------------------------------------------------------------
\newcommand{\PROVED}{\textsf{[proved here]}}
\newcommand{\STD}[1]{\textsf{[standard: #1]}}   % MUST carry a precise citation
\newcommand{\CONJ}{\textsf{[conjectural]}}
\newcommand{\COMP}{\textsf{[verified computationally]}}
\newcommand{\OPENSTAT}{\textsf{[open]}}

%----------------------------------------------------------------------
% Number systems, rings, groups
%----------------------------------------------------------------------
\newcommand{\Z}{\mathbb{Z}}
\newcommand{\Q}{\mathbb{Q}}
\newcommand{\Rl}{\mathbb{R}}              % the real numbers (NOT \R: see below)
\newcommand{\Cx}{\mathbb{C}}              % the complex numbers (NOT \C)
\newcommand{\F}{\mathbb{F}}
\newcommand{\Fl}{\F_\ell}
\newcommand{\Flf}{\F_{\ell^{f}}}          % residue field, residue degree f
\newcommand{\Hyp}{\mathcal{H}}            % the upper half plane
\newcommand{\GL}{\mathrm{GL}}
\newcommand{\SL}{\mathrm{SL}}
\newcommand{\Mat}{\mathrm{M}}             % matrix algebra: \Mat_2(\Z)
\newcommand{\Ga}{\Gamma}
\newcommand{\Gbdry}{\Ga_\infty}           % the boundary group <T> x <-I>
\newcommand{\De}{\Delta}                  % the monoid {g in M_2(Z) : det g > 0}
\newcommand{\Deprim}{\De^{\mathrm{prim}}}
\newcommand{\cring}{R}                    % the coefficient ring: F_ell, Z/ell^e, Z/m
\newcommand{\Rel}{\mathcal{R}}            % the relation set of the model (K,\Rel)

%----------------------------------------------------------------------
% Operators
%----------------------------------------------------------------------
\newcommand{\Hom}{\mathrm{Hom}}
\newcommand{\End}{\mathrm{End}}
\newcommand{\Ind}{\mathrm{Ind}}
\newcommand{\Sym}{\mathrm{Sym}}
\newcommand{\id}{\mathrm{id}}
\DeclareMathOperator{\res}{res}
\DeclareMathOperator{\cor}{cor}
\DeclareMathOperator{\conj}{conj}
\DeclareMathOperator{\adjg}{adj}
\DeclareMathOperator{\tr}{tr}
\DeclareMathOperator{\rank}{rank}
\DeclareMathOperator{\sgn}{sign}
\DeclareMathOperator{\fl}{fl}
\DeclareMathOperator{\im}{im}
\DeclareMathOperator{\coker}{coker}
\DeclareMathOperator{\soc}{soc}
\DeclareMathOperator{\rad}{rad}
\DeclareMathOperator{\Ext}{Ext}
\DeclareMathOperator{\Ann}{Ann}
\DeclareMathOperator{\MaxSpec}{MaxSpec}
\DeclareMathOperator{\Frob}{Frob}
\DeclareMathOperator{\supp}{supp}
\newcommand{\aug}{\varepsilon}            % augmentation C_0 -> Z

%----------------------------------------------------------------------
% The simplicial framework
%----------------------------------------------------------------------
\newcommand{\Ku}{K^{u}}                   % universal cover complex
\newcommand{\KR}{K/\Rel}                  % quotient complex
\newcommand{\prR}{\mathrm{pr}_{\Rel}}     % projection K -> K/\Rel
\newcommand{\pru}{\mathrm{pr}^{u}}        % projection K^u -> K/\Rel
\newcommand{\piOne}{\pi_1}
\newcommand{\LCS}{\mathbf{LCS}}
\newcommand{\Rep}{\mathbf{Rep}}
\newcommand{\LS}{\mathbf{LS}}
\newcommand{\cat}[1]{\mathbf{#1}}
\newcommand{\orient}[1]{o_{#1}}           % oriented q-simplices: \orient{q}(K)
% The distinguished lift of a simplex of K into the universal cover K^u, for the
% strict fundamental domain K \subseteq K^u of (H4)(iii).  The code spells this
% i(sigma); the document cannot, because i is sqrt(-1) (notation chapter, C15/C16).
\newcommand{\lift}[1]{\widetilde{#1}}
\newcommand{\Ceq}{C^{\bullet}_{\mathrm{eq}}}
\newcommand{\Cq}[1]{C^{#1}_{\mathrm{eq}}}
\newcommand{\Zq}[1]{Z^{#1}_{\mathrm{eq}}}
\newcommand{\Bq}[1]{B^{#1}_{\mathrm{eq}}}
\newcommand{\Heq}[1]{H^{#1}_{\mathrm{eq}}}

%----------------------------------------------------------------------
% Coefficients
%----------------------------------------------------------------------
\newcommand{\rhon}{\rho_n}                % the symmetric power representation
\newcommand{\rhodual}{\rho_n^{\vee}}      % the contragredient
\newcommand{\Vcoef}{V}                    % V := Sym^n(Z^2)
\newcommand{\VR}{V_{\cring}}              % V \otimes \cring
\newcommand{\Mn}{M_n}                     % Sym^n of F_ell^2 (synonym of \VR at R = F_ell)
\newcommand{\Dn}{D_n}                     % the divided-power / contragredient module
\newcommand{\cf}{f_{12}}                  % the normalised weight-12 cusp form (Ramanujan's Delta)

%----------------------------------------------------------------------
% The Hecke datum
%----------------------------------------------------------------------
\newcommand{\Th}{\Theta}                  % the fused cochain map Theta_g
\newcommand{\Thb}{\Theta^{\partial}}      % the boundary cochain map
\newcommand{\Xig}{\Xi_g}
\newcommand{\taub}{\tau_{\partial}}       % the boundary g-twisted chain datum
\newcommand{\Phig}{\Phi_g}                % the cone cochain map
\newcommand{\Amat}[1]{A^{(#1)}}           % cohomology-level Hecke matrix
\newcommand{\Acone}[1]{A_c^{(#1)}}
\newcommand{\Abdry}[1]{A_{\partial}^{(#1)}}
\newcommand{\Tdc}[1]{T_{\Ga #1 \Ga}}      % T_{Gamma g Gamma}
\newcommand{\diam}{T(p,p)}                % the diamond operator

%----------------------------------------------------------------------
% The pair, the cone, compact supports
%----------------------------------------------------------------------
\newcommand{\Xbulk}{X}                    % the Gamma-tessellation of the truncated half plane
\newcommand{\Xbdry}{\partial X}
\newcommand{\Kbdry}{K_{\partial}}         % the circle model 0-1-2-3
\newcommand{\rres}{r}                     % the restriction cochain map r = j^*
\newcommand{\Cone}[1]{C(#1)}              % the mapping cone
\newcommand{\Dcone}{D}                    % the cone differential
\newcommand{\coneinc}{\mathrm{inc}_c}     % C_eq(K_bdry) -> C(r)
\newcommand{\coneproj}{\mathrm{pr}_c}     % C(r) -> C_eq(K)[1]
\newcommand{\Hcpt}[1]{H^{#1}_{c}}
\newcommand{\etapair}{\eta}               % the pair homotopy
\newcommand{\hder}{h}                     % the derived cochain homotopy h = \tilde\eta
\newcommand{\Lam}{\Lambda}                % the unique C_0(bdry X) -> C_2(X) of eta-uniqueness

%----------------------------------------------------------------------
% The integral line (sec:integral, since [C] v1.8.0; notation
% sec:notation-integral)
%----------------------------------------------------------------------
\newcommand{\Iso}{\mathcal{I}}            % the isotropy-resolution model on X
\newcommand{\Isob}{\mathcal{I}^{\partial}} % its boundary subcomplex
\newcommand{\cmpI}{\mathrm{cmp}}          % the comparison map I_* -> C_*(X)
\newcommand{\CI}{C^{\bullet}_{\Iso}}      % Hom_{Z[Gamma]}(I_*, V_R)
\newcommand{\Hfr}[1]{H^{#1}_{\mathrm{fr}}} % H^q(Gamma;V) modulo torsion
\newcommand{\Zsix}{\widehat{\Z}_{6}}      % Z_2 x Z_3
\newcommand{\Qsix}{\widehat{\Q}_{6}}      % Q_2 x Q_3

%----------------------------------------------------------------------
% The Dickson line and the ordinary projector (sec:dicksontheta,
% sec:ordinary, since [C] v1.11.0-pre.1; notation sec:notation-dickson).
% Sans-serif L, D: the plain letters are taken (L the local system, D_n the
% contragredient module), rule A5.
%----------------------------------------------------------------------
\newcommand{\DkL}{\mathsf{L}}             % the first Dickson invariant x^ell y - x y^ell (the code's L)
\newcommand{\DkD}{\mathsf{D}}             % the second Dickson invariant, degree ell^2 - ell
\newcommand{\thL}{\theta_{\DkL}}          % multiplication by L on cochains and cohomology
\newcommand{\nplus}{n^{+}}               % n + ell + 1 (N is the Hecke level only)
\newcommand{\Qn}{Q_{n}}                   % the quotient Sym^{n+ell+1} / L Sym^n
\newcommand{\eord}{e_{\mathrm{ord}}}      % the ordinary projector of def:ordinary (e alone is lem:bulkperiodic's class)

%----------------------------------------------------------------------
% The involution
%----------------------------------------------------------------------
\newcommand{\inv}{\iota}                  % the involution diag(-1,1) / z -> -zbar
\newcommand{\invK}{\iota_K}               % its simplicial realisation on the 7-vertex model

%----------------------------------------------------------------------
% The Hecke algebra and its local pieces
%----------------------------------------------------------------------
\newcommand{\Talg}{\mathbb{T}}
\newcommand{\midl}{\mathfrak{m}}
\newcommand{\Htot}{H_{\mathrm{tot}}}

%----------------------------------------------------------------------
% Misc
%----------------------------------------------------------------------
\newcommand{\abs}[1]{\left| #1 \right|}
\newcommand{\smat}[4]{\begin{pmatrix} #1 & #2 \\ #3 & #4 \end{pmatrix}}
\newcommand{\code}[1]{\texttt{#1}}
\newcommand{\file}[1]{\texttt{#1}}

%----------------------------------------------------------------------
% Cross-document citation.  \HA{} survives for documents that still cite the
% OLD hecke_action.tex; inside the unified contract, use \ref{} instead.
%----------------------------------------------------------------------
\newcommand{\HA}[1]{[C,~\texttt{#1}]}
\newcommand{\vfive}[1]{[v5,~\texttt{#1}]}
\newcommand{\Espec}[1]{[E,~\texttt{#1}]}

%----------------------------------------------------------------------
% RETIRED SYMBOLS.  These macros existed in the pre-consolidation corpus with
% a meaning the notation chapter has retired.  Using one is an error, not a
% typographic choice.  (\K, \R, \C, \pR, \pu, \image were in the copied macro
% blocks of hecke_action.tex and monodromy_v5.tex.)
%----------------------------------------------------------------------
\newcommand{\gltwoRETIRED}[2]{%
  \PackageError{gl2_preamble}{Retired macro \string#1 -- #2}{%
    See the notation chapter (sec:notation) for the replacement.}}
\newcommand{\K}{\gltwoRETIRED{K}{ambiguous complex name; write K directly}}
\newcommand{\R}{\gltwoRETIRED{R}{R is the coefficient ring; write \string\Rl\space for the reals}}
\newcommand{\C}{\gltwoRETIRED{C}{write \string\Cx\space for the complex numbers}}
\newcommand{\pR}{\gltwoRETIRED{pR}{write \string\prR}}
\newcommand{\pu}{\gltwoRETIRED{pu}{write \string\pru}}
\newcommand{\image}{\gltwoRETIRED{image}{write \string\im}}
\newcommand{\M}{\gltwoRETIRED{M}{write \string\Mat\space for the matrix algebra}}

%
%   Every macro below is normative: the notation chapter of the contract
%   (gl2_notation.tex) is the prose form of this file, and the two must agree.
%   A symbol that is not declared here is not part of the fixed notation.
%
%   RETIRED MACROS.  Symbols retired by the notation chapter are redefined
%   here to raise a LaTeX error rather than silently typeset, so that text
%   re-homed from the old documents cannot carry an old meaning in unnoticed.
%   See the block "Retired symbols" at the end.
%======================================================================

\usepackage{amsmath, amssymb, amsthm}
\usepackage{mathtools}
\usepackage{geometry}
\usepackage{booktabs}
\usepackage{array}
\usepackage{longtable}
\usepackage{enumitem}
\usepackage{tikz}
\usepackage{tikz-cd}
\usepackage[colorlinks=true,linkcolor=blue,citecolor=blue,urlcolor=blue]{hyperref}
\setlist{nosep}
\emergencystretch=3em

%----------------------------------------------------------------------
% Theorem environments.  One counter, sectioned, so every statement has a
% unique number of the form <section>.<n> that can be cited from code.
%----------------------------------------------------------------------
\newtheorem{theorem}{Theorem}[section]
\newtheorem{lemma}[theorem]{Lemma}
\newtheorem{proposition}[theorem]{Proposition}
\newtheorem{corollary}[theorem]{Corollary}
\newtheorem{claim}[theorem]{Claim}

\theoremstyle{definition}
\newtheorem{definition}[theorem]{Definition}
\newtheorem{algorithm}[theorem]{Algorithm}
\newtheorem{obligation}[theorem]{Obligation}
\newtheorem{certificate}[theorem]{Certificate}

\theoremstyle{remark}
\newtheorem{remark}[theorem]{Remark}
\newtheorem{assumption}[theorem]{Assumption}
\newtheorem{example}[theorem]{Example}
\newtheorem{prediction}[theorem]{Prediction}
\newtheorem{openitem}[theorem]{Open construction item}

%----------------------------------------------------------------------
% Normative boxes.
%----------------------------------------------------------------------
\newenvironment{contractbox}
  {\par\medskip\noindent\begin{quote}\textbf{CONTRACT.}\ }
  {\end{quote}\medskip}
\newenvironment{decisionbox}
  {\par\medskip\noindent\begin{quote}\textbf{OWNER DECISION.}\ }
  {\end{quote}\medskip}
\newenvironment{erratumbox}
  {\par\medskip\noindent\begin{quote}\textbf{ERRATUM.}\ }
  {\end{quote}\medskip}

%----------------------------------------------------------------------
% STATUS tags.  Every non-definitional assertion carries exactly one.
%----------------------------------------------------------------------
\newcommand{\PROVED}{\textsf{[proved here]}}
\newcommand{\STD}[1]{\textsf{[standard: #1]}}   % MUST carry a precise citation
\newcommand{\CONJ}{\textsf{[conjectural]}}
\newcommand{\COMP}{\textsf{[verified computationally]}}
\newcommand{\OPENSTAT}{\textsf{[open]}}

%----------------------------------------------------------------------
% Number systems, rings, groups
%----------------------------------------------------------------------
\newcommand{\Z}{\mathbb{Z}}
\newcommand{\Q}{\mathbb{Q}}
\newcommand{\Rl}{\mathbb{R}}              % the real numbers (NOT \R: see below)
\newcommand{\Cx}{\mathbb{C}}              % the complex numbers (NOT \C)
\newcommand{\F}{\mathbb{F}}
\newcommand{\Fl}{\F_\ell}
\newcommand{\Flf}{\F_{\ell^{f}}}          % residue field, residue degree f
\newcommand{\Hyp}{\mathcal{H}}            % the upper half plane
\newcommand{\GL}{\mathrm{GL}}
\newcommand{\SL}{\mathrm{SL}}
\newcommand{\Mat}{\mathrm{M}}             % matrix algebra: \Mat_2(\Z)
\newcommand{\Ga}{\Gamma}
\newcommand{\Gbdry}{\Ga_\infty}           % the boundary group <T> x <-I>
\newcommand{\De}{\Delta}                  % the monoid {g in M_2(Z) : det g > 0}
\newcommand{\Deprim}{\De^{\mathrm{prim}}}
\newcommand{\cring}{R}                    % the coefficient ring: F_ell, Z/ell^e, Z/m
\newcommand{\Rel}{\mathcal{R}}            % the relation set of the model (K,\Rel)

%----------------------------------------------------------------------
% Operators
%----------------------------------------------------------------------
\newcommand{\Hom}{\mathrm{Hom}}
\newcommand{\End}{\mathrm{End}}
\newcommand{\Ind}{\mathrm{Ind}}
\newcommand{\Sym}{\mathrm{Sym}}
\newcommand{\id}{\mathrm{id}}
\DeclareMathOperator{\res}{res}
\DeclareMathOperator{\cor}{cor}
\DeclareMathOperator{\conj}{conj}
\DeclareMathOperator{\adjg}{adj}
\DeclareMathOperator{\tr}{tr}
\DeclareMathOperator{\rank}{rank}
\DeclareMathOperator{\sgn}{sign}
\DeclareMathOperator{\fl}{fl}
\DeclareMathOperator{\im}{im}
\DeclareMathOperator{\coker}{coker}
\DeclareMathOperator{\soc}{soc}
\DeclareMathOperator{\rad}{rad}
\DeclareMathOperator{\Ext}{Ext}
\DeclareMathOperator{\Ann}{Ann}
\DeclareMathOperator{\MaxSpec}{MaxSpec}
\DeclareMathOperator{\Frob}{Frob}
\DeclareMathOperator{\supp}{supp}
\newcommand{\aug}{\varepsilon}            % augmentation C_0 -> Z

%----------------------------------------------------------------------
% The simplicial framework
%----------------------------------------------------------------------
\newcommand{\Ku}{K^{u}}                   % universal cover complex
\newcommand{\KR}{K/\Rel}                  % quotient complex
\newcommand{\prR}{\mathrm{pr}_{\Rel}}     % projection K -> K/\Rel
\newcommand{\pru}{\mathrm{pr}^{u}}        % projection K^u -> K/\Rel
\newcommand{\piOne}{\pi_1}
\newcommand{\LCS}{\mathbf{LCS}}
\newcommand{\Rep}{\mathbf{Rep}}
\newcommand{\LS}{\mathbf{LS}}
\newcommand{\cat}[1]{\mathbf{#1}}
\newcommand{\orient}[1]{o_{#1}}           % oriented q-simplices: \orient{q}(K)
% The distinguished lift of a simplex of K into the universal cover K^u, for the
% strict fundamental domain K \subseteq K^u of (H4)(iii).  The code spells this
% i(sigma); the document cannot, because i is sqrt(-1) (notation chapter, C15/C16).
\newcommand{\lift}[1]{\widetilde{#1}}
\newcommand{\Ceq}{C^{\bullet}_{\mathrm{eq}}}
\newcommand{\Cq}[1]{C^{#1}_{\mathrm{eq}}}
\newcommand{\Zq}[1]{Z^{#1}_{\mathrm{eq}}}
\newcommand{\Bq}[1]{B^{#1}_{\mathrm{eq}}}
\newcommand{\Heq}[1]{H^{#1}_{\mathrm{eq}}}

%----------------------------------------------------------------------
% Coefficients
%----------------------------------------------------------------------
\newcommand{\rhon}{\rho_n}                % the symmetric power representation
\newcommand{\rhodual}{\rho_n^{\vee}}      % the contragredient
\newcommand{\Vcoef}{V}                    % V := Sym^n(Z^2)
\newcommand{\VR}{V_{\cring}}              % V \otimes \cring
\newcommand{\Mn}{M_n}                     % Sym^n of F_ell^2 (synonym of \VR at R = F_ell)
\newcommand{\Dn}{D_n}                     % the divided-power / contragredient module
\newcommand{\cf}{f_{12}}                  % the normalised weight-12 cusp form (Ramanujan's Delta)

%----------------------------------------------------------------------
% The Hecke datum
%----------------------------------------------------------------------
\newcommand{\Th}{\Theta}                  % the fused cochain map Theta_g
\newcommand{\Thb}{\Theta^{\partial}}      % the boundary cochain map
\newcommand{\Xig}{\Xi_g}
\newcommand{\taub}{\tau_{\partial}}       % the boundary g-twisted chain datum
\newcommand{\Phig}{\Phi_g}                % the cone cochain map
\newcommand{\Amat}[1]{A^{(#1)}}           % cohomology-level Hecke matrix
\newcommand{\Acone}[1]{A_c^{(#1)}}
\newcommand{\Abdry}[1]{A_{\partial}^{(#1)}}
\newcommand{\Tdc}[1]{T_{\Ga #1 \Ga}}      % T_{Gamma g Gamma}
\newcommand{\diam}{T(p,p)}                % the diamond operator

%----------------------------------------------------------------------
% The pair, the cone, compact supports
%----------------------------------------------------------------------
\newcommand{\Xbulk}{X}                    % the Gamma-tessellation of the truncated half plane
\newcommand{\Xbdry}{\partial X}
\newcommand{\Kbdry}{K_{\partial}}         % the circle model 0-1-2-3
\newcommand{\rres}{r}                     % the restriction cochain map r = j^*
\newcommand{\Cone}[1]{C(#1)}              % the mapping cone
\newcommand{\Dcone}{D}                    % the cone differential
\newcommand{\coneinc}{\mathrm{inc}_c}     % C_eq(K_bdry) -> C(r)
\newcommand{\coneproj}{\mathrm{pr}_c}     % C(r) -> C_eq(K)[1]
\newcommand{\Hcpt}[1]{H^{#1}_{c}}
\newcommand{\etapair}{\eta}               % the pair homotopy
\newcommand{\hder}{h}                     % the derived cochain homotopy h = \tilde\eta
\newcommand{\Lam}{\Lambda}                % the unique C_0(bdry X) -> C_2(X) of eta-uniqueness

%----------------------------------------------------------------------
% The integral line (sec:integral, since [C] v1.8.0; notation
% sec:notation-integral)
%----------------------------------------------------------------------
\newcommand{\Iso}{\mathcal{I}}            % the isotropy-resolution model on X
\newcommand{\Isob}{\mathcal{I}^{\partial}} % its boundary subcomplex
\newcommand{\cmpI}{\mathrm{cmp}}          % the comparison map I_* -> C_*(X)
\newcommand{\CI}{C^{\bullet}_{\Iso}}      % Hom_{Z[Gamma]}(I_*, V_R)
\newcommand{\Hfr}[1]{H^{#1}_{\mathrm{fr}}} % H^q(Gamma;V) modulo torsion
\newcommand{\Zsix}{\widehat{\Z}_{6}}      % Z_2 x Z_3
\newcommand{\Qsix}{\widehat{\Q}_{6}}      % Q_2 x Q_3

%----------------------------------------------------------------------
% The Dickson line and the ordinary projector (sec:dicksontheta,
% sec:ordinary, since [C] v1.11.0-pre.1; notation sec:notation-dickson).
% Sans-serif L, D: the plain letters are taken (L the local system, D_n the
% contragredient module), rule A5.
%----------------------------------------------------------------------
\newcommand{\DkL}{\mathsf{L}}             % the first Dickson invariant x^ell y - x y^ell (the code's L)
\newcommand{\DkD}{\mathsf{D}}             % the second Dickson invariant, degree ell^2 - ell
\newcommand{\thL}{\theta_{\DkL}}          % multiplication by L on cochains and cohomology
\newcommand{\nplus}{n^{+}}               % n + ell + 1 (N is the Hecke level only)
\newcommand{\Qn}{Q_{n}}                   % the quotient Sym^{n+ell+1} / L Sym^n
\newcommand{\eord}{e_{\mathrm{ord}}}      % the ordinary projector of def:ordinary (e alone is lem:bulkperiodic's class)

%----------------------------------------------------------------------
% The involution
%----------------------------------------------------------------------
\newcommand{\inv}{\iota}                  % the involution diag(-1,1) / z -> -zbar
\newcommand{\invK}{\iota_K}               % its simplicial realisation on the 7-vertex model

%----------------------------------------------------------------------
% The Hecke algebra and its local pieces
%----------------------------------------------------------------------
\newcommand{\Talg}{\mathbb{T}}
\newcommand{\midl}{\mathfrak{m}}
\newcommand{\Htot}{H_{\mathrm{tot}}}

%----------------------------------------------------------------------
% Misc
%----------------------------------------------------------------------
\newcommand{\abs}[1]{\left| #1 \right|}
\newcommand{\smat}[4]{\begin{pmatrix} #1 & #2 \\ #3 & #4 \end{pmatrix}}
\newcommand{\code}[1]{\texttt{#1}}
\newcommand{\file}[1]{\texttt{#1}}

%----------------------------------------------------------------------
% Cross-document citation.  \HA{} survives for documents that still cite the
% OLD hecke_action.tex; inside the unified contract, use \ref{} instead.
%----------------------------------------------------------------------
\newcommand{\HA}[1]{[C,~\texttt{#1}]}
\newcommand{\vfive}[1]{[v5,~\texttt{#1}]}
\newcommand{\Espec}[1]{[E,~\texttt{#1}]}

%----------------------------------------------------------------------
% RETIRED SYMBOLS.  These macros existed in the pre-consolidation corpus with
% a meaning the notation chapter has retired.  Using one is an error, not a
% typographic choice.  (\K, \R, \C, \pR, \pu, \image were in the copied macro
% blocks of hecke_action.tex and monodromy_v5.tex.)
%----------------------------------------------------------------------
\newcommand{\gltwoRETIRED}[2]{%
  \PackageError{gl2_preamble}{Retired macro \string#1 -- #2}{%
    See the notation chapter (sec:notation) for the replacement.}}
\newcommand{\K}{\gltwoRETIRED{K}{ambiguous complex name; write K directly}}
\newcommand{\R}{\gltwoRETIRED{R}{R is the coefficient ring; write \string\Rl\space for the reals}}
\newcommand{\C}{\gltwoRETIRED{C}{write \string\Cx\space for the complex numbers}}
\newcommand{\pR}{\gltwoRETIRED{pR}{write \string\prR}}
\newcommand{\pu}{\gltwoRETIRED{pu}{write \string\pru}}
\newcommand{\image}{\gltwoRETIRED{image}{write \string\im}}
\newcommand{\M}{\gltwoRETIRED{M}{write \string\Mat\space for the matrix algebra}}

%
%   Every macro below is normative: the notation chapter of the contract
%   (gl2_notation.tex) is the prose form of this file, and the two must agree.
%   A symbol that is not declared here is not part of the fixed notation.
%
%   RETIRED MACROS.  Symbols retired by the notation chapter are redefined
%   here to raise a LaTeX error rather than silently typeset, so that text
%   re-homed from the old documents cannot carry an old meaning in unnoticed.
%   See the block "Retired symbols" at the end.
%======================================================================

\usepackage{amsmath, amssymb, amsthm}
\usepackage{mathtools}
\usepackage{geometry}
\usepackage{booktabs}
\usepackage{array}
\usepackage{longtable}
\usepackage{enumitem}
\usepackage{tikz}
\usepackage{tikz-cd}
\usepackage[colorlinks=true,linkcolor=blue,citecolor=blue,urlcolor=blue]{hyperref}
\setlist{nosep}
\emergencystretch=3em

%----------------------------------------------------------------------
% Theorem environments.  One counter, sectioned, so every statement has a
% unique number of the form <section>.<n> that can be cited from code.
%----------------------------------------------------------------------
\newtheorem{theorem}{Theorem}[section]
\newtheorem{lemma}[theorem]{Lemma}
\newtheorem{proposition}[theorem]{Proposition}
\newtheorem{corollary}[theorem]{Corollary}
\newtheorem{claim}[theorem]{Claim}

\theoremstyle{definition}
\newtheorem{definition}[theorem]{Definition}
\newtheorem{algorithm}[theorem]{Algorithm}
\newtheorem{obligation}[theorem]{Obligation}
\newtheorem{certificate}[theorem]{Certificate}

\theoremstyle{remark}
\newtheorem{remark}[theorem]{Remark}
\newtheorem{assumption}[theorem]{Assumption}
\newtheorem{example}[theorem]{Example}
\newtheorem{prediction}[theorem]{Prediction}
\newtheorem{openitem}[theorem]{Open construction item}

%----------------------------------------------------------------------
% Normative boxes.
%----------------------------------------------------------------------
\newenvironment{contractbox}
  {\par\medskip\noindent\begin{quote}\textbf{CONTRACT.}\ }
  {\end{quote}\medskip}
\newenvironment{decisionbox}
  {\par\medskip\noindent\begin{quote}\textbf{OWNER DECISION.}\ }
  {\end{quote}\medskip}
\newenvironment{erratumbox}
  {\par\medskip\noindent\begin{quote}\textbf{ERRATUM.}\ }
  {\end{quote}\medskip}

%----------------------------------------------------------------------
% STATUS tags.  Every non-definitional assertion carries exactly one.
%----------------------------------------------------------------------
\newcommand{\PROVED}{\textsf{[proved here]}}
\newcommand{\STD}[1]{\textsf{[standard: #1]}}   % MUST carry a precise citation
\newcommand{\CONJ}{\textsf{[conjectural]}}
\newcommand{\COMP}{\textsf{[verified computationally]}}
\newcommand{\OPENSTAT}{\textsf{[open]}}

%----------------------------------------------------------------------
% Number systems, rings, groups
%----------------------------------------------------------------------
\newcommand{\Z}{\mathbb{Z}}
\newcommand{\Q}{\mathbb{Q}}
\newcommand{\Rl}{\mathbb{R}}              % the real numbers (NOT \R: see below)
\newcommand{\Cx}{\mathbb{C}}              % the complex numbers (NOT \C)
\newcommand{\F}{\mathbb{F}}
\newcommand{\Fl}{\F_\ell}
\newcommand{\Flf}{\F_{\ell^{f}}}          % residue field, residue degree f
\newcommand{\Hyp}{\mathcal{H}}            % the upper half plane
\newcommand{\GL}{\mathrm{GL}}
\newcommand{\SL}{\mathrm{SL}}
\newcommand{\Mat}{\mathrm{M}}             % matrix algebra: \Mat_2(\Z)
\newcommand{\Ga}{\Gamma}
\newcommand{\Gbdry}{\Ga_\infty}           % the boundary group <T> x <-I>
\newcommand{\De}{\Delta}                  % the monoid {g in M_2(Z) : det g > 0}
\newcommand{\Deprim}{\De^{\mathrm{prim}}}
\newcommand{\cring}{R}                    % the coefficient ring: F_ell, Z/ell^e, Z/m
\newcommand{\Rel}{\mathcal{R}}            % the relation set of the model (K,\Rel)

%----------------------------------------------------------------------
% Operators
%----------------------------------------------------------------------
\newcommand{\Hom}{\mathrm{Hom}}
\newcommand{\End}{\mathrm{End}}
\newcommand{\Ind}{\mathrm{Ind}}
\newcommand{\Sym}{\mathrm{Sym}}
\newcommand{\id}{\mathrm{id}}
\DeclareMathOperator{\res}{res}
\DeclareMathOperator{\cor}{cor}
\DeclareMathOperator{\conj}{conj}
\DeclareMathOperator{\adjg}{adj}
\DeclareMathOperator{\tr}{tr}
\DeclareMathOperator{\rank}{rank}
\DeclareMathOperator{\sgn}{sign}
\DeclareMathOperator{\fl}{fl}
\DeclareMathOperator{\im}{im}
\DeclareMathOperator{\coker}{coker}
\DeclareMathOperator{\soc}{soc}
\DeclareMathOperator{\rad}{rad}
\DeclareMathOperator{\Ext}{Ext}
\DeclareMathOperator{\Ann}{Ann}
\DeclareMathOperator{\MaxSpec}{MaxSpec}
\DeclareMathOperator{\Frob}{Frob}
\DeclareMathOperator{\supp}{supp}
\newcommand{\aug}{\varepsilon}            % augmentation C_0 -> Z

%----------------------------------------------------------------------
% The simplicial framework
%----------------------------------------------------------------------
\newcommand{\Ku}{K^{u}}                   % universal cover complex
\newcommand{\KR}{K/\Rel}                  % quotient complex
\newcommand{\prR}{\mathrm{pr}_{\Rel}}     % projection K -> K/\Rel
\newcommand{\pru}{\mathrm{pr}^{u}}        % projection K^u -> K/\Rel
\newcommand{\piOne}{\pi_1}
\newcommand{\LCS}{\mathbf{LCS}}
\newcommand{\Rep}{\mathbf{Rep}}
\newcommand{\LS}{\mathbf{LS}}
\newcommand{\cat}[1]{\mathbf{#1}}
\newcommand{\orient}[1]{o_{#1}}           % oriented q-simplices: \orient{q}(K)
% The distinguished lift of a simplex of K into the universal cover K^u, for the
% strict fundamental domain K \subseteq K^u of (H4)(iii).  The code spells this
% i(sigma); the document cannot, because i is sqrt(-1) (notation chapter, C15/C16).
\newcommand{\lift}[1]{\widetilde{#1}}
\newcommand{\Ceq}{C^{\bullet}_{\mathrm{eq}}}
\newcommand{\Cq}[1]{C^{#1}_{\mathrm{eq}}}
\newcommand{\Zq}[1]{Z^{#1}_{\mathrm{eq}}}
\newcommand{\Bq}[1]{B^{#1}_{\mathrm{eq}}}
\newcommand{\Heq}[1]{H^{#1}_{\mathrm{eq}}}

%----------------------------------------------------------------------
% Coefficients
%----------------------------------------------------------------------
\newcommand{\rhon}{\rho_n}                % the symmetric power representation
\newcommand{\rhodual}{\rho_n^{\vee}}      % the contragredient
\newcommand{\Vcoef}{V}                    % V := Sym^n(Z^2)
\newcommand{\VR}{V_{\cring}}              % V \otimes \cring
\newcommand{\Mn}{M_n}                     % Sym^n of F_ell^2 (synonym of \VR at R = F_ell)
\newcommand{\Dn}{D_n}                     % the divided-power / contragredient module
\newcommand{\cf}{f_{12}}                  % the normalised weight-12 cusp form (Ramanujan's Delta)

%----------------------------------------------------------------------
% The Hecke datum
%----------------------------------------------------------------------
\newcommand{\Th}{\Theta}                  % the fused cochain map Theta_g
\newcommand{\Thb}{\Theta^{\partial}}      % the boundary cochain map
\newcommand{\Xig}{\Xi_g}
\newcommand{\taub}{\tau_{\partial}}       % the boundary g-twisted chain datum
\newcommand{\Phig}{\Phi_g}                % the cone cochain map
\newcommand{\Amat}[1]{A^{(#1)}}           % cohomology-level Hecke matrix
\newcommand{\Acone}[1]{A_c^{(#1)}}
\newcommand{\Abdry}[1]{A_{\partial}^{(#1)}}
\newcommand{\Tdc}[1]{T_{\Ga #1 \Ga}}      % T_{Gamma g Gamma}
\newcommand{\diam}{T(p,p)}                % the diamond operator

%----------------------------------------------------------------------
% The pair, the cone, compact supports
%----------------------------------------------------------------------
\newcommand{\Xbulk}{X}                    % the Gamma-tessellation of the truncated half plane
\newcommand{\Xbdry}{\partial X}
\newcommand{\Kbdry}{K_{\partial}}         % the circle model 0-1-2-3
\newcommand{\rres}{r}                     % the restriction cochain map r = j^*
\newcommand{\Cone}[1]{C(#1)}              % the mapping cone
\newcommand{\Dcone}{D}                    % the cone differential
\newcommand{\coneinc}{\mathrm{inc}_c}     % C_eq(K_bdry) -> C(r)
\newcommand{\coneproj}{\mathrm{pr}_c}     % C(r) -> C_eq(K)[1]
\newcommand{\Hcpt}[1]{H^{#1}_{c}}
\newcommand{\etapair}{\eta}               % the pair homotopy
\newcommand{\hder}{h}                     % the derived cochain homotopy h = \tilde\eta
\newcommand{\Lam}{\Lambda}                % the unique C_0(bdry X) -> C_2(X) of eta-uniqueness

%----------------------------------------------------------------------
% The integral line (sec:integral, since [C] v1.8.0; notation
% sec:notation-integral)
%----------------------------------------------------------------------
\newcommand{\Iso}{\mathcal{I}}            % the isotropy-resolution model on X
\newcommand{\Isob}{\mathcal{I}^{\partial}} % its boundary subcomplex
\newcommand{\cmpI}{\mathrm{cmp}}          % the comparison map I_* -> C_*(X)
\newcommand{\CI}{C^{\bullet}_{\Iso}}      % Hom_{Z[Gamma]}(I_*, V_R)
\newcommand{\Hfr}[1]{H^{#1}_{\mathrm{fr}}} % H^q(Gamma;V) modulo torsion
\newcommand{\Zsix}{\widehat{\Z}_{6}}      % Z_2 x Z_3
\newcommand{\Qsix}{\widehat{\Q}_{6}}      % Q_2 x Q_3

%----------------------------------------------------------------------
% The Dickson line and the ordinary projector (sec:dicksontheta,
% sec:ordinary, since [C] v1.11.0-pre.1; notation sec:notation-dickson).
% Sans-serif L, D: the plain letters are taken (L the local system, D_n the
% contragredient module), rule A5.
%----------------------------------------------------------------------
\newcommand{\DkL}{\mathsf{L}}             % the first Dickson invariant x^ell y - x y^ell (the code's L)
\newcommand{\DkD}{\mathsf{D}}             % the second Dickson invariant, degree ell^2 - ell
\newcommand{\thL}{\theta_{\DkL}}          % multiplication by L on cochains and cohomology
\newcommand{\nplus}{n^{+}}               % n + ell + 1 (N is the Hecke level only)
\newcommand{\Qn}{Q_{n}}                   % the quotient Sym^{n+ell+1} / L Sym^n
\newcommand{\eord}{e_{\mathrm{ord}}}      % the ordinary projector of def:ordinary (e alone is lem:bulkperiodic's class)

%----------------------------------------------------------------------
% The involution
%----------------------------------------------------------------------
\newcommand{\inv}{\iota}                  % the involution diag(-1,1) / z -> -zbar
\newcommand{\invK}{\iota_K}               % its simplicial realisation on the 7-vertex model

%----------------------------------------------------------------------
% The Hecke algebra and its local pieces
%----------------------------------------------------------------------
\newcommand{\Talg}{\mathbb{T}}
\newcommand{\midl}{\mathfrak{m}}
\newcommand{\Htot}{H_{\mathrm{tot}}}

%----------------------------------------------------------------------
% Misc
%----------------------------------------------------------------------
\newcommand{\abs}[1]{\left| #1 \right|}
\newcommand{\smat}[4]{\begin{pmatrix} #1 & #2 \\ #3 & #4 \end{pmatrix}}
\newcommand{\code}[1]{\texttt{#1}}
\newcommand{\file}[1]{\texttt{#1}}

%----------------------------------------------------------------------
% Cross-document citation.  \HA{} survives for documents that still cite the
% OLD hecke_action.tex; inside the unified contract, use \ref{} instead.
%----------------------------------------------------------------------
\newcommand{\HA}[1]{[C,~\texttt{#1}]}
\newcommand{\vfive}[1]{[v5,~\texttt{#1}]}
\newcommand{\Espec}[1]{[E,~\texttt{#1}]}

%----------------------------------------------------------------------
% RETIRED SYMBOLS.  These macros existed in the pre-consolidation corpus with
% a meaning the notation chapter has retired.  Using one is an error, not a
% typographic choice.  (\K, \R, \C, \pR, \pu, \image were in the copied macro
% blocks of hecke_action.tex and monodromy_v5.tex.)
%----------------------------------------------------------------------
\newcommand{\gltwoRETIRED}[2]{%
  \PackageError{gl2_preamble}{Retired macro \string#1 -- #2}{%
    See the notation chapter (sec:notation) for the replacement.}}
\newcommand{\K}{\gltwoRETIRED{K}{ambiguous complex name; write K directly}}
\newcommand{\R}{\gltwoRETIRED{R}{R is the coefficient ring; write \string\Rl\space for the reals}}
\newcommand{\C}{\gltwoRETIRED{C}{write \string\Cx\space for the complex numbers}}
\newcommand{\pR}{\gltwoRETIRED{pR}{write \string\prR}}
\newcommand{\pu}{\gltwoRETIRED{pu}{write \string\pru}}
\newcommand{\image}{\gltwoRETIRED{image}{write \string\im}}
\newcommand{\M}{\gltwoRETIRED{M}{write \string\Mat\space for the matrix algebra}}

%
%   Every macro below is normative: the notation chapter of the contract
%   (gl2_notation.tex) is the prose form of this file, and the two must agree.
%   A symbol that is not declared here is not part of the fixed notation.
%
%   RETIRED MACROS.  Symbols retired by the notation chapter are redefined
%   here to raise a LaTeX error rather than silently typeset, so that text
%   re-homed from the old documents cannot carry an old meaning in unnoticed.
%   See the block "Retired symbols" at the end.
%======================================================================

\usepackage{amsmath, amssymb, amsthm}
\usepackage{mathtools}
\usepackage{geometry}
\usepackage{booktabs}
\usepackage{array}
\usepackage{longtable}
\usepackage{enumitem}
\usepackage{tikz}
\usepackage{tikz-cd}
\usepackage[colorlinks=true,linkcolor=blue,citecolor=blue,urlcolor=blue]{hyperref}
\setlist{nosep}
\emergencystretch=3em

%----------------------------------------------------------------------
% Theorem environments.  One counter, sectioned, so every statement has a
% unique number of the form <section>.<n> that can be cited from code.
%----------------------------------------------------------------------
\newtheorem{theorem}{Theorem}[section]
\newtheorem{lemma}[theorem]{Lemma}
\newtheorem{proposition}[theorem]{Proposition}
\newtheorem{corollary}[theorem]{Corollary}
\newtheorem{claim}[theorem]{Claim}

\theoremstyle{definition}
\newtheorem{definition}[theorem]{Definition}
\newtheorem{algorithm}[theorem]{Algorithm}
\newtheorem{obligation}[theorem]{Obligation}
\newtheorem{certificate}[theorem]{Certificate}

\theoremstyle{remark}
\newtheorem{remark}[theorem]{Remark}
\newtheorem{assumption}[theorem]{Assumption}
\newtheorem{example}[theorem]{Example}
\newtheorem{prediction}[theorem]{Prediction}
\newtheorem{openitem}[theorem]{Open construction item}

%----------------------------------------------------------------------
% Normative boxes.
%----------------------------------------------------------------------
\newenvironment{contractbox}
  {\par\medskip\noindent\begin{quote}\textbf{CONTRACT.}\ }
  {\end{quote}\medskip}
\newenvironment{decisionbox}
  {\par\medskip\noindent\begin{quote}\textbf{OWNER DECISION.}\ }
  {\end{quote}\medskip}
\newenvironment{erratumbox}
  {\par\medskip\noindent\begin{quote}\textbf{ERRATUM.}\ }
  {\end{quote}\medskip}

%----------------------------------------------------------------------
% STATUS tags.  Every non-definitional assertion carries exactly one.
%----------------------------------------------------------------------
\newcommand{\PROVED}{\textsf{[proved here]}}
\newcommand{\STD}[1]{\textsf{[standard: #1]}}   % MUST carry a precise citation
\newcommand{\CONJ}{\textsf{[conjectural]}}
\newcommand{\COMP}{\textsf{[verified computationally]}}
\newcommand{\OPENSTAT}{\textsf{[open]}}

%----------------------------------------------------------------------
% Number systems, rings, groups
%----------------------------------------------------------------------
\newcommand{\Z}{\mathbb{Z}}
\newcommand{\Q}{\mathbb{Q}}
\newcommand{\Rl}{\mathbb{R}}              % the real numbers (NOT \R: see below)
\newcommand{\Cx}{\mathbb{C}}              % the complex numbers (NOT \C)
\newcommand{\F}{\mathbb{F}}
\newcommand{\Fl}{\F_\ell}
\newcommand{\Flf}{\F_{\ell^{f}}}          % residue field, residue degree f
\newcommand{\Hyp}{\mathcal{H}}            % the upper half plane
\newcommand{\GL}{\mathrm{GL}}
\newcommand{\SL}{\mathrm{SL}}
\newcommand{\Mat}{\mathrm{M}}             % matrix algebra: \Mat_2(\Z)
\newcommand{\Ga}{\Gamma}
\newcommand{\Gbdry}{\Ga_\infty}           % the boundary group <T> x <-I>
\newcommand{\De}{\Delta}                  % the monoid {g in M_2(Z) : det g > 0}
\newcommand{\Deprim}{\De^{\mathrm{prim}}}
\newcommand{\cring}{R}                    % the coefficient ring: F_ell, Z/ell^e, Z/m
\newcommand{\Rel}{\mathcal{R}}            % the relation set of the model (K,\Rel)

%----------------------------------------------------------------------
% Operators
%----------------------------------------------------------------------
\newcommand{\Hom}{\mathrm{Hom}}
\newcommand{\End}{\mathrm{End}}
\newcommand{\Ind}{\mathrm{Ind}}
\newcommand{\Sym}{\mathrm{Sym}}
\newcommand{\id}{\mathrm{id}}
\DeclareMathOperator{\res}{res}
\DeclareMathOperator{\cor}{cor}
\DeclareMathOperator{\conj}{conj}
\DeclareMathOperator{\adjg}{adj}
\DeclareMathOperator{\tr}{tr}
\DeclareMathOperator{\rank}{rank}
\DeclareMathOperator{\sgn}{sign}
\DeclareMathOperator{\fl}{fl}
\DeclareMathOperator{\im}{im}
\DeclareMathOperator{\coker}{coker}
\DeclareMathOperator{\soc}{soc}
\DeclareMathOperator{\rad}{rad}
\DeclareMathOperator{\Ext}{Ext}
\DeclareMathOperator{\Ann}{Ann}
\DeclareMathOperator{\MaxSpec}{MaxSpec}
\DeclareMathOperator{\Frob}{Frob}
\DeclareMathOperator{\supp}{supp}
\newcommand{\aug}{\varepsilon}            % augmentation C_0 -> Z

%----------------------------------------------------------------------
% The simplicial framework
%----------------------------------------------------------------------
\newcommand{\Ku}{K^{u}}                   % universal cover complex
\newcommand{\KR}{K/\Rel}                  % quotient complex
\newcommand{\prR}{\mathrm{pr}_{\Rel}}     % projection K -> K/\Rel
\newcommand{\pru}{\mathrm{pr}^{u}}        % projection K^u -> K/\Rel
\newcommand{\piOne}{\pi_1}
\newcommand{\LCS}{\mathbf{LCS}}
\newcommand{\Rep}{\mathbf{Rep}}
\newcommand{\LS}{\mathbf{LS}}
\newcommand{\cat}[1]{\mathbf{#1}}
\newcommand{\orient}[1]{o_{#1}}           % oriented q-simplices: \orient{q}(K)
% The distinguished lift of a simplex of K into the universal cover K^u, for the
% strict fundamental domain K \subseteq K^u of (H4)(iii).  The code spells this
% i(sigma); the document cannot, because i is sqrt(-1) (notation chapter, C15/C16).
\newcommand{\lift}[1]{\widetilde{#1}}
\newcommand{\Ceq}{C^{\bullet}_{\mathrm{eq}}}
\newcommand{\Cq}[1]{C^{#1}_{\mathrm{eq}}}
\newcommand{\Zq}[1]{Z^{#1}_{\mathrm{eq}}}
\newcommand{\Bq}[1]{B^{#1}_{\mathrm{eq}}}
\newcommand{\Heq}[1]{H^{#1}_{\mathrm{eq}}}

%----------------------------------------------------------------------
% Coefficients
%----------------------------------------------------------------------
\newcommand{\rhon}{\rho_n}                % the symmetric power representation
\newcommand{\rhodual}{\rho_n^{\vee}}      % the contragredient
\newcommand{\Vcoef}{V}                    % V := Sym^n(Z^2)
\newcommand{\VR}{V_{\cring}}              % V \otimes \cring
\newcommand{\Mn}{M_n}                     % Sym^n of F_ell^2 (synonym of \VR at R = F_ell)
\newcommand{\Dn}{D_n}                     % the divided-power / contragredient module
\newcommand{\cf}{f_{12}}                  % the normalised weight-12 cusp form (Ramanujan's Delta)

%----------------------------------------------------------------------
% The Hecke datum
%----------------------------------------------------------------------
\newcommand{\Th}{\Theta}                  % the fused cochain map Theta_g
\newcommand{\Thb}{\Theta^{\partial}}      % the boundary cochain map
\newcommand{\Xig}{\Xi_g}
\newcommand{\taub}{\tau_{\partial}}       % the boundary g-twisted chain datum
\newcommand{\Phig}{\Phi_g}                % the cone cochain map
\newcommand{\Amat}[1]{A^{(#1)}}           % cohomology-level Hecke matrix
\newcommand{\Acone}[1]{A_c^{(#1)}}
\newcommand{\Abdry}[1]{A_{\partial}^{(#1)}}
\newcommand{\Tdc}[1]{T_{\Ga #1 \Ga}}      % T_{Gamma g Gamma}
\newcommand{\diam}{T(p,p)}                % the diamond operator

%----------------------------------------------------------------------
% The pair, the cone, compact supports
%----------------------------------------------------------------------
\newcommand{\Xbulk}{X}                    % the Gamma-tessellation of the truncated half plane
\newcommand{\Xbdry}{\partial X}
\newcommand{\Kbdry}{K_{\partial}}         % the circle model 0-1-2-3
\newcommand{\rres}{r}                     % the restriction cochain map r = j^*
\newcommand{\Cone}[1]{C(#1)}              % the mapping cone
\newcommand{\Dcone}{D}                    % the cone differential
\newcommand{\coneinc}{\mathrm{inc}_c}     % C_eq(K_bdry) -> C(r)
\newcommand{\coneproj}{\mathrm{pr}_c}     % C(r) -> C_eq(K)[1]
\newcommand{\Hcpt}[1]{H^{#1}_{c}}
\newcommand{\etapair}{\eta}               % the pair homotopy
\newcommand{\hder}{h}                     % the derived cochain homotopy h = \tilde\eta
\newcommand{\Lam}{\Lambda}                % the unique C_0(bdry X) -> C_2(X) of eta-uniqueness

%----------------------------------------------------------------------
% The integral line (sec:integral, since [C] v1.8.0; notation
% sec:notation-integral)
%----------------------------------------------------------------------
\newcommand{\Iso}{\mathcal{I}}            % the isotropy-resolution model on X
\newcommand{\Isob}{\mathcal{I}^{\partial}} % its boundary subcomplex
\newcommand{\cmpI}{\mathrm{cmp}}          % the comparison map I_* -> C_*(X)
\newcommand{\CI}{C^{\bullet}_{\Iso}}      % Hom_{Z[Gamma]}(I_*, V_R)
\newcommand{\Hfr}[1]{H^{#1}_{\mathrm{fr}}} % H^q(Gamma;V) modulo torsion
\newcommand{\Zsix}{\widehat{\Z}_{6}}      % Z_2 x Z_3
\newcommand{\Qsix}{\widehat{\Q}_{6}}      % Q_2 x Q_3

%----------------------------------------------------------------------
% The Dickson line and the ordinary projector (sec:dicksontheta,
% sec:ordinary, since [C] v1.11.0-pre.1; notation sec:notation-dickson).
% Sans-serif L, D: the plain letters are taken (L the local system, D_n the
% contragredient module), rule A5.
%----------------------------------------------------------------------
\newcommand{\DkL}{\mathsf{L}}             % the first Dickson invariant x^ell y - x y^ell (the code's L)
\newcommand{\DkD}{\mathsf{D}}             % the second Dickson invariant, degree ell^2 - ell
\newcommand{\thL}{\theta_{\DkL}}          % multiplication by L on cochains and cohomology
\newcommand{\nplus}{n^{+}}               % n + ell + 1 (N is the Hecke level only)
\newcommand{\Qn}{Q_{n}}                   % the quotient Sym^{n+ell+1} / L Sym^n
\newcommand{\eord}{e_{\mathrm{ord}}}      % the ordinary projector of def:ordinary (e alone is lem:bulkperiodic's class)

%----------------------------------------------------------------------
% The involution
%----------------------------------------------------------------------
\newcommand{\inv}{\iota}                  % the involution diag(-1,1) / z -> -zbar
\newcommand{\invK}{\iota_K}               % its simplicial realisation on the 7-vertex model

%----------------------------------------------------------------------
% The Hecke algebra and its local pieces
%----------------------------------------------------------------------
\newcommand{\Talg}{\mathbb{T}}
\newcommand{\midl}{\mathfrak{m}}
\newcommand{\Htot}{H_{\mathrm{tot}}}

%----------------------------------------------------------------------
% Misc
%----------------------------------------------------------------------
\newcommand{\abs}[1]{\left| #1 \right|}
\newcommand{\smat}[4]{\begin{pmatrix} #1 & #2 \\ #3 & #4 \end{pmatrix}}
\newcommand{\code}[1]{\texttt{#1}}
\newcommand{\file}[1]{\texttt{#1}}

%----------------------------------------------------------------------
% Cross-document citation.  \HA{} survives for documents that still cite the
% OLD hecke_action.tex; inside the unified contract, use \ref{} instead.
%----------------------------------------------------------------------
\newcommand{\HA}[1]{[C,~\texttt{#1}]}
\newcommand{\vfive}[1]{[v5,~\texttt{#1}]}
\newcommand{\Espec}[1]{[E,~\texttt{#1}]}

%----------------------------------------------------------------------
% RETIRED SYMBOLS.  These macros existed in the pre-consolidation corpus with
% a meaning the notation chapter has retired.  Using one is an error, not a
% typographic choice.  (\K, \R, \C, \pR, \pu, \image were in the copied macro
% blocks of hecke_action.tex and monodromy_v5.tex.)
%----------------------------------------------------------------------
\newcommand{\gltwoRETIRED}[2]{%
  \PackageError{gl2_preamble}{Retired macro \string#1 -- #2}{%
    See the notation chapter (sec:notation) for the replacement.}}
\newcommand{\K}{\gltwoRETIRED{K}{ambiguous complex name; write K directly}}
\newcommand{\R}{\gltwoRETIRED{R}{R is the coefficient ring; write \string\Rl\space for the reals}}
\newcommand{\C}{\gltwoRETIRED{C}{write \string\Cx\space for the complex numbers}}
\newcommand{\pR}{\gltwoRETIRED{pR}{write \string\prR}}
\newcommand{\pu}{\gltwoRETIRED{pu}{write \string\pru}}
\newcommand{\image}{\gltwoRETIRED{image}{write \string\im}}
\newcommand{\M}{\gltwoRETIRED{M}{write \string\Mat\space for the matrix algebra}}
